
        \documentclass[fleqn]{article}      
        \usepackage{amsmath}
        \usepackage{fontspec}
        \usepackage{kotex} % 한국어 지원  

        \begin{document}      
        쉬움: \\\\
1. 다음 중 적분의 정의는 무엇인가요? \\\\
   a) 미분의 역작용 \\\\
   b) 함수의 곡선 아래 면적 \\\\
   c) 함수의 기울기 \\\\
   d) 함수의 최대값 \\\\

2. ∫2x dx의 결과는 무엇인가요? \\\\
   a) x^2 + C \\\\
   b) 2x^2 + C \\\\
   c) 2x + C \\\\
   d) x^2 + 2C \\\\

3. ∫(3x^2)dx의 결과는 무엇인가요? \\\\
   a) x^3 + C \\\\
   b) 3/3 x^3 + C \\\\
   c) 3x^3 + C \\\\
   d) 3/4 x^4 + C \\\\

4. 적분을 통해 구할 수 있는 값은 무엇인가요? \\\\
   a) 함수의 최대값 \\\\
   b) 함수의 기울기 \\\\
   c) 곡선 아래의 면적 \\\\
   d) 함수의 평균값 \\\\

5. ∫(1/x)dx의 결과는 무엇인가요? \\\\
   a) ln|x| + C \\\\
   b) x + C \\\\
   c) 1/x + C \\\\
   d) x^2 + C \\\\

중간: \\\\
1. ∫(3x^2 + 4x)dx를 계산하시오. \\\\
   
2. ∫(sin x)dx의 결과는 무엇인가요? \\\\
   
3. ∫(2x^3 - 5x + 1)dx를 계산하시오. \\\\
   
4. ∫(e^x)dx의 결과는 무엇인가요? \\\\
   
5. ∫(x^2 * e^x)dx를 부분적분을 이용하여 계산하시오. \\\\

어려움: \\\\
1. ∫(x^3 * ln(x))dx를 부분적분을 이용하여 계산하시오. \\\\
   
2. ∫(1/(x^2 + 1))dx의 결과는 무엇인가요? \\\\
   
3. ∫(x^2 * cos(x))dx를 부분적분을 이용하여 계산하시오. \\\\
   
4. ∫(x^4 * e^x)dx를 부분적분을 이용하여 계산하시오. \\\\
   
5. ∫(sin^2(x))dx의 결과를 구하시오. \\\\

      
        \end{document} 
    